% =====================================================================
%  [개정본 2026-07-27] 5장 논의와 한계 (한국어)
%  원본: sections/05-ko.tex (무변경 보존)
%  개정: '본 연구가 주장하지 않는 것' -> '해석 범위와 식별 조건'.
%        6개 항목의 정보량은 유지하고 문법만 방법론으로 전환한다.
% =====================================================================

\section{논의와 한계}
\label{sec:discussion}

SPR-IRL이 복원한 것은 세 유형의 행동을 하나의 공통 척도 위에 나란히 놓고, 그 중
어디까지가 선호이고 어디부터가 시장 청산의 산술인지 구분한 그림이다. 모멘텀 축에서
외국인과 개인이 정반대로 갈리는 대비는 자기상관 통제 후에도 살아남으며, 이 중
외국인 측은 무선호 합성 대역 밖이므로 추세추종 선호로 읽힌다. 개인 측 계수와 컨텍스트
주효과는 청산 제약이 유도하는 범위 안에 있어 집계 수급의 구조로 읽는다. 기관은 시차
교차연관을 제외하면 식별되는 선호를 보이지 않으며 이는 합산 범주의 상쇄로 설명된다.
이 결과들은 세 유형에 동일한 특징과 동일한 모형 형태를 부여한 뒤 파라미터만 분리
학습해 얻은 것이므로 설계가 심어둔 대비가 아니다.

\subsection{해석 범위와 식별 조건}
\label{sec:scope}

프로토콜의 마지막 단계로, 각 결과가 어떤 조건에서 식별되고 어디까지 해석되는지를
명시한다.

\textbf{(1) 수익률 예측가능성.} 본 연구의 지표는 수급의 방향과 강도에 대한 재현
오차이며 수익률, 샤프지수, 거래비용을 포함하지 않는다. 표~\ref{tab:performance}에서
어떤 매매 규칙도 도출되지 않는다.

\textbf{(2) 수급 예측과 선호 복원의 분리.} (F1)--(F3) 사양은 지속성 기준선에 방향
정확도와 상관계수에서 밀린다(\S\ref{sec:oos}). 이는 결함이 아니라 특징 집합의 설계가
정의하는 과제의 차이다. 선호 복원은 해석 가능한 공유 특징을, 수급 예측은 자기지연을
요구한다. 지연 자기행동을 추가한 승격 사양은 기준선을 넘어서면서($0.652/0.574/
0.629$) 모멘텀 대비와 $\alpha^{FX}$ 대비를 보존하고 $\alpha^{KOSPI}$를
반전시킨다(표~\ref{tab:spec34}). 두 과제를 같은 모형으로 겸하려 할 때 어떤 해석이
살아남는지에 대한 실측이다.

\textbf{(3) 인과효과.} 복원된 가중치는 표준화된 관측 특징 사이의 조건부 연관이다.
특징--컨텍스트 상관이 $-0.29$에서 $0.22$ 범위이므로 개별 계수의 크기는 다른 항의
통제 조건에 의존한다.

\textbf{(4) 유형 간 음의 연관의 행동적 해석.} 시차 교차연관 계수와 개인 측 거울상
계수는 시장 청산 제약이 유도하는 성분과 분리되지 않으며, 그 유도 성분의 크기는
\S\ref{sec:beta}에서 정량화했다(관측의 84--214\%). 이것은 집계 수급 자료에서 유형 간
계수를 해석할 때의 식별 경계이며, 본 논문은 그 경계 안에서---기계 대역을 초과하는
외국인 계수만 선호로, 나머지는 집계 구조로---결과를 서술한다.

\textbf{(5) 처분효과의 검정.} 개인의 손실구간 부호는 문헌과 정합하나 월 블록
부트스트랩을 통과하지 못했고 자기지연 통제에서 더 약해지며(77.8\%), 평균단가 대용치는
계좌별 실현손익이 아니라 집계 거래로 만든 근사치다. 따라서 처분효과의 검정이 아니라
부호 정합의 관측이다.

\textbf{(6) 종목 간 일반화.} 단일 종목 설정은 유형 간 이질성 분리를 위한 통제이며
다른 종목과 시장으로의 확장은 검증되지 않았다.

\subsection{향후 과제}
\label{sec:future}

세 방향이 검정 가능한 형태로 남는다. 첫째, 기관 세부주체 분해. \S\ref{sec:null}의
null이 상쇄에서 온다면---연기금ㆍ투신ㆍ보험ㆍ사모ㆍ금융투자가 서로 다른 목적함수로
거래해 합산 순매수가 0 근방에 머문다는 가설---KRX 투자자별 거래실적의 세부 주체
자료로 직접 검정할 수 있다. 본 논문의 대칭 추정과 V1--V4는 세부 주체 수에 관계없이
그대로 적용되며, 예측은 명확하다. 분해 후 각 주체에서 기관 합산에서 사라졌던 계수가
반대 부호로 나타나야 한다. 둘째, 추가 종목 일반화. 청산 유도 성분의 크기는 종목별
투자자 구성비 $\mathbb{E}[W_i/W_j]$와 흡수 몫의 함수이므로, 외국인 지분과 유동성이
다른 종목에서는 \S\ref{sec:beta}의 유도 성분 대비 관측 계수의 비율이 예측 가능하게
달라져야 한다. 이것이 거울상의 기계적 성분과 선호 성분을 종목 횡단면에서 분리하는
설계다. 셋째, 근시안 가정의 해제. 명제~\ref{prop:lasso}가 보이듯 한 걸음만 내다보는
설정에서 보상은 정책과 겹치지만, 상태 동역학과 다기간 가치를 도입하면 두 대상이
갈라지고 회귀로는 도달할 수 없는 구조적 복원이 된다. 그 확장에서도 SPR-IRL의 대칭
추정과 검증 프로토콜은 그대로 쓰인다.
